% !TEX spellcheck = en_US



% ~~~~~~~~~~~~~~~~~~~~~~~~~~~~~~~~~~~~~~~ V
% (NC) 2026 Haluk Bingol
% github.com/halukbingol/LaTeX-Templates
%
% Licensees may copy, distribute, display, and perform the work and make derivative works 
% and remixes based on it only for non-commercial purposes. 
% ~~~~~~~~~~~~~~~~~~~~~~~~~~~~~~~~~~~~~~~ A




% =====================================================================
%  Strings in Java -- Freshman Level: First Steps with Text
%  ---------------------------------------------------------------------
%  Built on hbCisLaTeXTemplate.tex with the libraries in ../../hblib/.
%  Programs and their outputs are read from codeoutput/:
%     codeoutput/<Name>.java   the program
%     codeoutput/<Name>.in     keyboard input (optional)
%     codeoutput/<Name>.txt    what it printed, made by:  sh run_examples.sh
% =====================================================================

\documentclass[aspectratio=169,10pt,compress]{beamer}

% ~~~~~~~~~~~~~~~~~~~~~~~~~~~~~~~~~~~~~~~~~~~~~~~~~~~~~~~~~~~~~~~~~ V hblib
	\usepackage{../../hblib/hblibPresentation} % lib presentation: theme, i/N, \hSection, algorithm2e, ...
	\usepackage{../../hblib/hblibCodeoutput} % lib code + output: \lstcode, \lstcodedown, ...
	\usepackage{../../hblib/hblibDatastructures} % lib data structures: \hString, \hStack, \hQueue
% ~~~~~~~~~~~~~~~~~~~~~~~~~~~~~~~~~~~~~~~~~~~~~~~~~~~~~~~~~~~~~~~~~ A hblib


% xcolor, array (and the L{width} column), listings, tikz and algorithm2e come from the hblib libraries
\usepackage[T1]{fontenc}
\usepackage{lmodern}
\usepackage{booktabs}

\definecolor{gitred}{RGB}{240,80,50}
\newcommand{\cmd}[1]{\texttt{\color{gitred}#1}}

% ---------------------------------------------------------------------
\title{Strings in Java}
\subtitle{Freshman Level: First Steps with Text}
\author[Bingol]{Haluk O. Bingol}

\institute[FCIS]{
    Faculty of Computer and Information Sciences\\
    Yeditepe University
}
\date{
	{\tiny \hbTimeStamp}
}

\begin{document}




% ~~~~~~~~~~~~~~~~~~~~~~~~~~~~~~~~~~~~~~ V frm
\begin{frame}[t,fragile]
  \titlepage
\end{frame}
% ~~~~~~~~~~~~~~~~~~~~~~~~~~~~~~~~~~~~~~ A frm




% ~~~~~~~~~~~~~~~~~~~~~~~~~~~~~~~~~~~~~~ V frm
\begin{frame}[t,fragile] \frametitle{Outline}
  \begin{columns}[T]
    \begin{column}{0.48\textwidth}
      \tableofcontents[sections={1-4}, hideallsubsections]
    \end{column}
    \begin{column}{0.48\textwidth}
      \tableofcontents[sections={5-8}, hideallsubsections]
    \end{column}
  \end{columns}
\end{frame}
% ~~~~~~~~~~~~~~~~~~~~~~~~~~~~~~~~~~~~~~ A frm

% ~~~~~~~~~~~~~~~~~~~~~~~~~~~~~~~~~~~~~~ V
\hSection[Setup]{Running Java on Your Computer}{}{}
% ~~~~~~~~~~~~~~~~~~~~~~~~~~~~~~~~~~~~~~ A

% ~~~~~~~~~~~~~~~~~~~~~~~~~~~~~~~~~~~~~~ V
\subsection[Install]{Install and check}
% ~~~~~~~~~~~~~~~~~~~~~~~~~~~~~~~~~~~~~~ A




% ~~~~~~~~~~~~~~~~~~~~~~~~~~~~~~~~~~~~~~ V frm
\begin{frame}[t,fragile] \frametitle{Install Java and check it}
  \small You need a \textbf{JDK} (Java Development Kit), for example Eclipse Temurin 21 from \texttt{adoptium.net}.
  \begin{columns}[T]
    \begin{column}{0.48\textwidth}
      \begin{block}{Windows}
        \footnotesize
        \begin{enumerate}
          \item Download the \texttt{.msi} installer, run it, and tick \emph{Set JAVA\_HOME} and \emph{Add to PATH}
          \item Open \textbf{Command Prompt} (\texttt{Win} key, type \texttt{cmd})
          \item Check:
        \end{enumerate}


% +++++++++++++++++++++++++++++++++++++++ V lst
%: -Lst.
\begin{lstlisting}[style=codesnippet, language=bash]
java -version
javac -version
\end{lstlisting}
% +++++++++++++++++++++++++++++++++++++++ A lst



      \end{block}
    \end{column}
    \begin{column}{0.48\textwidth}
      \begin{block}{macOS}
        \footnotesize
        \begin{enumerate}
          \item Download the \texttt{.pkg} installer (Apple silicon: \emph{aarch64}; Intel: \emph{x64}) and run it
          \item Open \textbf{Terminal} (\texttt{Cmd}+\texttt{Space}, type \texttt{Terminal})
          \item Check:
        \end{enumerate}


% +++++++++++++++++++++++++++++++++++++++ V lst
%: -Lst.
\begin{lstlisting}[style=codesnippet, language=bash]
java -version
javac -version
\end{lstlisting}
% +++++++++++++++++++++++++++++++++++++++ A lst



      \end{block}
    \end{column}
  \end{columns}
  \small Both commands should print a version number such as \texttt{21.0.10}.
\end{frame}
% ~~~~~~~~~~~~~~~~~~~~~~~~~~~~~~~~~~~~~~ A frm

% ~~~~~~~~~~~~~~~~~~~~~~~~~~~~~~~~~~~~~~ V
\subsection[Run]{Compile and run}
% ~~~~~~~~~~~~~~~~~~~~~~~~~~~~~~~~~~~~~~ A




% ~~~~~~~~~~~~~~~~~~~~~~~~~~~~~~~~~~~~~~ V frm
\begin{frame}[t,fragile] \frametitle{Compile and run a program}
  \begin{columns}[T]
    \begin{column}{0.48\textwidth}
      \begin{block}{Windows (Command Prompt)}


% +++++++++++++++++++++++++++++++++++++++ V lst
%: -Lst.
\begin{lstlisting}[style=codesnippet, language=bash]
cd C:\Users\ayse\cse101
chcp 65001
javac -encoding UTF-8 Greeting.java
java Greeting
\end{lstlisting}
% +++++++++++++++++++++++++++++++++++++++ A lst



        \footnotesize \cmd{chcp 65001} lets the console show Turkish letters (ç, ğ, ı, ş).
      \end{block}
    \end{column}
    \begin{column}{0.48\textwidth}
      \begin{block}{macOS (Terminal)}


% +++++++++++++++++++++++++++++++++++++++ V lst
%: -Lst.
\begin{lstlisting}[style=codesnippet, language=bash]
cd ~/cse101
javac Greeting.java
java Greeting
\end{lstlisting}
% +++++++++++++++++++++++++++++++++++++++ A lst



        \footnotesize Terminal is UTF-8 already; no extra step.
      \end{block}
    \end{column}
  \end{columns}
  \vspace{0.3em}
  \begin{itemize}\small
    \item \cmd{javac} \textbf{compiles} \texttt{Greeting.java} into \texttt{Greeting.class}; \cmd{java} \textbf{runs} it (no \texttt{.class} at the end)
    \item Shortcut since Java 11 for one-file programs: \cmd{java Greeting.java}
    \item The file name must match the class name: \texttt{public class Greeting} lives in \texttt{Greeting.java}
  \end{itemize}
\end{frame}
% ~~~~~~~~~~~~~~~~~~~~~~~~~~~~~~~~~~~~~~ A frm

% ~~~~~~~~~~~~~~~~~~~~~~~~~~~~~~~~~~~~~~ V
\hSection[First]{Your First Strings}{}{}
% ~~~~~~~~~~~~~~~~~~~~~~~~~~~~~~~~~~~~~~ A

% ~~~~~~~~~~~~~~~~~~~~~~~~~~~~~~~~~~~~~~ V
\subsection[What]{What is a string?}
% ~~~~~~~~~~~~~~~~~~~~~~~~~~~~~~~~~~~~~~ A




% ~~~~~~~~~~~~~~~~~~~~~~~~~~~~~~~~~~~~~~ V frm
\begin{frame}[t,fragile] \frametitle{What is a string?}
  \begin{columns}[T]
    \begin{column}{0.5\textwidth}
      A \textbf{string} is a piece of text: a sequence of characters between double quotes.


% +++++++++++++++++++++++++++++++++++++++ V lst
%: -Lst.
\begin{lstlisting}[style=codesnippet, language=Java]
String city = "İstanbul";
String empty = "";
char letter = 'A';      // one char
\end{lstlisting}
% +++++++++++++++++++++++++++++++++++++++ A lst



      \begin{itemize}\small
        \item \texttt{String} (capital S) is the type
        \item Double quotes \texttt{"..."} for strings, single quotes \texttt{'.'} for one \texttt{char}
        \item \texttt{+} glues strings together (\textbf{concatenation})
      \end{itemize}
    \end{column}
    \begin{column}{0.46\textwidth}
      \small Every character has a position, starting at \textbf{0}:
      \vspace{0.6em}
      \begin{center}
        \begin{tikzpicture}\hString[0.48]{İ,s,t,a,n,b,u,l}\end{tikzpicture}
      \end{center}
      \vspace{0.4em}
      \small \texttt{city.length()} is 8,\\
      \texttt{city.charAt(0)} is \texttt{'İ'},\\
      \texttt{city.charAt(7)} is \texttt{'l'}.
    \end{column}
  \end{columns}
\end{frame}
% ~~~~~~~~~~~~~~~~~~~~~~~~~~~~~~~~~~~~~~ A frm

% ~~~~~~~~~~~~~~~~~~~~~~~~~~~~~~~~~~~~~~ V
\subsection[Plus]{Joining with +}
% ~~~~~~~~~~~~~~~~~~~~~~~~~~~~~~~~~~~~~~ A




% ~~~~~~~~~~~~~~~~~~~~~~~~~~~~~~~~~~~~~~ V frm
\begin{frame}[t,fragile] \frametitle{Joining strings with \texttt{+}}
  \lstcoderightstep[0.7][basicstyle=\ttfamily\scriptsize][basicstyle=\ttfamily\scriptsize]{codeoutput/Greeting.java}{Java}{codeoutput/Greeting.txt}

  \vspace{0.3em}
  \small \hRemark{Look at the last line:} \texttt{"Class of " + year} is a string, then \texttt{+ 4} appends ``4''.
  Use parentheses to add numbers first.
\end{frame}
% ~~~~~~~~~~~~~~~~~~~~~~~~~~~~~~~~~~~~~~ A frm




% ~~~~~~~~~~~~~~~~~~~~~~~~~~~~~~~~~~~~~~ V frm
\begin{frame}[t,fragile] \frametitle{Puzzle: what does \texttt{+} do here?}
  \lstcoderightstep[0.7][basicstyle=\ttfamily\scriptsize][basicstyle=\ttfamily\scriptsize]{codeoutput/PlusPuzzle.java}{Java}{codeoutput/PlusPuzzle.txt}

  \vspace{0.3em}
  \small \texttt{+} works \textbf{left to right}. Number $+$ number adds; as soon as one side is a \texttt{String}, it joins.
  Two \texttt{char}s are numbers: \texttt{'a' + 'b'} is $97 + 98$.
\end{frame}
% ~~~~~~~~~~~~~~~~~~~~~~~~~~~~~~~~~~~~~~ A frm

% ~~~~~~~~~~~~~~~~~~~~~~~~~~~~~~~~~~~~~~ V
\hSection[Input]{Reading Text}{}{}
% ~~~~~~~~~~~~~~~~~~~~~~~~~~~~~~~~~~~~~~ A

% ~~~~~~~~~~~~~~~~~~~~~~~~~~~~~~~~~~~~~~ V
\subsection[Scanner]{Reading with Scanner}
% ~~~~~~~~~~~~~~~~~~~~~~~~~~~~~~~~~~~~~~ A




% ~~~~~~~~~~~~~~~~~~~~~~~~~~~~~~~~~~~~~~ V frm
\begin{frame}[t,fragile] \frametitle{Reading a line with \texttt{Scanner}}
  \lstcoderightstep[0.7][basicstyle=\ttfamily\scriptsize][basicstyle=\ttfamily\scriptsize]{codeoutput/ReadName.java}{Java}{codeoutput/ReadName.txt}

  \vspace{0.3em}
  \small Here the input comes from the file \texttt{ReadName.in} (\texttt{Ayşe Yılmaz}, \texttt{İstanbul}),
  so the typed text is not shown after the prompts. \texttt{nextLine()} reads a whole line, spaces included.
\end{frame}
% ~~~~~~~~~~~~~~~~~~~~~~~~~~~~~~~~~~~~~~ A frm

% ~~~~~~~~~~~~~~~~~~~~~~~~~~~~~~~~~~~~~~ V
\subsection[Trap]{The nextLine trap}
% ~~~~~~~~~~~~~~~~~~~~~~~~~~~~~~~~~~~~~~ A




% ~~~~~~~~~~~~~~~~~~~~~~~~~~~~~~~~~~~~~~ V frm
\begin{frame}[t,fragile] \frametitle{The \texttt{nextInt} / \texttt{nextLine} trap}
  \lstcoderightstep[0.7][basicstyle=\ttfamily\scriptsize][basicstyle=\ttfamily\scriptsize]{codeoutput/NextLineTrap.java}{Java}{codeoutput/NextLineTrap.txt}

  \vspace{0.3em}
  \small Input: \texttt{19}, Enter, \texttt{Ali Can}, Enter. \texttt{nextInt()} stops before the Enter, so the first
  \texttt{nextLine()} returns the empty rest of that line. \hRemark{Fix:} call \texttt{in.nextLine()} once after \texttt{nextInt()}.
\end{frame}
% ~~~~~~~~~~~~~~~~~~~~~~~~~~~~~~~~~~~~~~ A frm

% ~~~~~~~~~~~~~~~~~~~~~~~~~~~~~~~~~~~~~~ V
\hSection[Loop]{Character by Character}{}{}
% ~~~~~~~~~~~~~~~~~~~~~~~~~~~~~~~~~~~~~~ A

% ~~~~~~~~~~~~~~~~~~~~~~~~~~~~~~~~~~~~~~ V
\subsection[Loop]{Looping over a string}
% ~~~~~~~~~~~~~~~~~~~~~~~~~~~~~~~~~~~~~~ A




% ~~~~~~~~~~~~~~~~~~~~~~~~~~~~~~~~~~~~~~ V frm
\begin{frame}[t,fragile] \frametitle{Looping over a string}
  \begin{columns}[T]
    \begin{column}{0.56\textwidth}


% +++++++++++++++++++++++++++++++++++++++ V lst
%: -Lst.
\begin{lstlisting}[style=codesnippet, language=Java]
for (int i = 0; i < s.length(); i++) {
    char c = s.charAt(i);
    // use c
}
\end{lstlisting}
% +++++++++++++++++++++++++++++++++++++++ A lst



      \begin{itemize}\small
        \item Start at \texttt{0}
        \item Stop \textbf{before} \texttt{s.length()}: use \texttt{<}, not \texttt{<=}
        \item \texttt{charAt(i)} gives the \texttt{char} at position \texttt{i}
      \end{itemize}
    \end{column}
    \begin{column}{0.4\textwidth}
      \begin{center}
      \begin{tikzpicture}
        \hString[0.6]{J,a,v,a}
        \foreach \i in {0,1,2,3} \draw[-{Stealth}, thick, orange] (\i*0.6+0.3,-0.55) -- (\i*0.6+0.3,-0.05);
        \node[font=\scriptsize, orange] at (1.2,-0.8) {\texttt{i} = 0, 1, 2, 3};
        \node[font=\scriptsize] at (1.2,-1.25) {\texttt{i} = 4 = \texttt{length()}: stop};
      \end{tikzpicture}
      \end{center}
    \end{column}
  \end{columns}
\end{frame}
% ~~~~~~~~~~~~~~~~~~~~~~~~~~~~~~~~~~~~~~ A frm




% ~~~~~~~~~~~~~~~~~~~~~~~~~~~~~~~~~~~~~~ V frm
\begin{frame}[t,fragile] \frametitle{Spelling a word}
  \lstcoderightstep[0.7][basicstyle=\ttfamily\scriptsize][basicstyle=\ttfamily\scriptsize]{codeoutput/Spell.java}{Java}{codeoutput/Spell.txt}
\end{frame}
% ~~~~~~~~~~~~~~~~~~~~~~~~~~~~~~~~~~~~~~ A frm

% ~~~~~~~~~~~~~~~~~~~~~~~~~~~~~~~~~~~~~~ V
\subsection[Count]{Counting}
% ~~~~~~~~~~~~~~~~~~~~~~~~~~~~~~~~~~~~~~ A




% ~~~~~~~~~~~~~~~~~~~~~~~~~~~~~~~~~~~~~~ V frm
\begin{frame}[t,fragile] \frametitle{Counting vowels}
  \lstcoderightstep[0.7][basicstyle=\ttfamily\scriptsize][basicstyle=\ttfamily\scriptsize]{codeoutput/CountVowels.java}{Java}{codeoutput/CountVowels.txt}

  \vspace{0.3em}
  \small A handy trick: \texttt{vowels.indexOf(c) >= 0} means ``\texttt{c} is one of the characters in \texttt{vowels}''.
\end{frame}
% ~~~~~~~~~~~~~~~~~~~~~~~~~~~~~~~~~~~~~~ A frm

% ~~~~~~~~~~~~~~~~~~~~~~~~~~~~~~~~~~~~~~ V
\hSection[Compare]{Comparing Strings}{}{}
% ~~~~~~~~~~~~~~~~~~~~~~~~~~~~~~~~~~~~~~ A

% ~~~~~~~~~~~~~~~~~~~~~~~~~~~~~~~~~~~~~~ V
\subsection[Equals]{equals, not ==}
% ~~~~~~~~~~~~~~~~~~~~~~~~~~~~~~~~~~~~~~ A




% ~~~~~~~~~~~~~~~~~~~~~~~~~~~~~~~~~~~~~~ V frm
\begin{frame}[t,fragile] \frametitle{The most famous string bug}
  \lstcoderightstep[0.7][basicstyle=\ttfamily\scriptsize][basicstyle=\ttfamily\scriptsize]{codeoutput/EqualsTrap.java}{Java}{codeoutput/EqualsTrap.txt}

  \vspace{0.3em}
  \small The user typed \texttt{yes}, yet \texttt{answer == "yes"} is \texttt{false}: \texttt{==} asks whether they are
  the \textbf{same object}. \hRemark{Always compare text with} \texttt{equals}: \texttt{answer.equals("yes")}.
\end{frame}
% ~~~~~~~~~~~~~~~~~~~~~~~~~~~~~~~~~~~~~~ A frm




% ~~~~~~~~~~~~~~~~~~~~~~~~~~~~~~~~~~~~~~ V frm
\begin{frame}[t,fragile] \frametitle{Comparing strings: the rules}
  \small
  \begin{tabular}{@{}L{0.4\textwidth}L{0.5\textwidth}@{}}
    \toprule
    \textbf{Question}                     & \textbf{Write} \\
    \midrule
    same text?                            & \texttt{a.equals(b)} \\
    same text, ignoring case?             & \texttt{a.equalsIgnoreCase(b)} \\
    which comes first?                    & \texttt{a.compareTo(b) < 0} \\
    is it empty?                          & \texttt{a.isEmpty()} \\
    does it contain \texttt{"x"}?         & \texttt{a.contains("x")} \\
    does it start with \texttt{"Dr."}?    & \texttt{a.startsWith("Dr.")} \\
    \bottomrule
  \end{tabular}

  \vspace{0.8em}
  \begin{alertblock}{Never}
    \texttt{if (a == b)} with strings. It may even work in a small test (with literals), then fail with user input.
  \end{alertblock}
\end{frame}
% ~~~~~~~~~~~~~~~~~~~~~~~~~~~~~~~~~~~~~~ A frm

% ~~~~~~~~~~~~~~~~~~~~~~~~~~~~~~~~~~~~~~ V
\hSection[Errors]{When Things Go Wrong}{}{}
% ~~~~~~~~~~~~~~~~~~~~~~~~~~~~~~~~~~~~~~ A

% ~~~~~~~~~~~~~~~~~~~~~~~~~~~~~~~~~~~~~~ V
\subsection[Index]{Index out of bounds}
% ~~~~~~~~~~~~~~~~~~~~~~~~~~~~~~~~~~~~~~ A




% ~~~~~~~~~~~~~~~~~~~~~~~~~~~~~~~~~~~~~~ V frm
\begin{frame}[t,fragile] \frametitle{Index out of bounds}
  \lstcodedownstep[basicstyle=\ttfamily\scriptsize][basicstyle=\ttfamily\scriptsize, linerange={1-2,13-13}]{codeoutput/OutOfBounds.java}{Java}{codeoutput/OutOfBounds.txt}

  \small \hRemark{Read the message:} the exception name, \texttt{Index 4 out of bounds for length 4}, and the
  line of \emph{your} file (\texttt{OutOfBounds.java:5}). Valid indices are \texttt{0..length()-1}.
  (The 10 lines from inside Java, \texttt{at java.base/...}, are skipped here.)
\end{frame}
% ~~~~~~~~~~~~~~~~~~~~~~~~~~~~~~~~~~~~~~ A frm

% ~~~~~~~~~~~~~~~~~~~~~~~~~~~~~~~~~~~~~~ V
\subsection[Null]{null}
% ~~~~~~~~~~~~~~~~~~~~~~~~~~~~~~~~~~~~~~ A




% ~~~~~~~~~~~~~~~~~~~~~~~~~~~~~~~~~~~~~~ V frm
\begin{frame}[t,fragile] \frametitle{A string that is not there: \texttt{null}}
  \lstcodedownstep[basicstyle=\ttfamily\scriptsize][basicstyle=\ttfamily\scriptsize]{codeoutput/NullName.java}{Java}{codeoutput/NullName.txt}

  \small \texttt{name} refers to no object. Printing it shows \texttt{null}, but calling a method on it crashes.
  Java even tells us which variable was \texttt{null}.
\end{frame}
% ~~~~~~~~~~~~~~~~~~~~~~~~~~~~~~~~~~~~~~ A frm

% ~~~~~~~~~~~~~~~~~~~~~~~~~~~~~~~~~~~~~~ V
\hSection[Problems]{Small Problems}{}{}
% ~~~~~~~~~~~~~~~~~~~~~~~~~~~~~~~~~~~~~~ A

% ~~~~~~~~~~~~~~~~~~~~~~~~~~~~~~~~~~~~~~ V
\subsection[Reverse]{Reversing}
% ~~~~~~~~~~~~~~~~~~~~~~~~~~~~~~~~~~~~~~ A




% ~~~~~~~~~~~~~~~~~~~~~~~~~~~~~~~~~~~~~~ V frm
\begin{frame}[t,fragile] \frametitle{Reversing a string}
  \lstcoderightstep[0.7][basicstyle=\ttfamily\scriptsize][basicstyle=\ttfamily\scriptsize]{codeoutput/Reverse.java}{Java}{codeoutput/Reverse.txt}

  \vspace{0.3em}
  \small The loop walks from the last index down to 0. \texttt{StringBuilder} does the same in one call, and faster.
\end{frame}
% ~~~~~~~~~~~~~~~~~~~~~~~~~~~~~~~~~~~~~~ A frm

% ~~~~~~~~~~~~~~~~~~~~~~~~~~~~~~~~~~~~~~ V
\subsection[Palindrome]{Palindromes}
% ~~~~~~~~~~~~~~~~~~~~~~~~~~~~~~~~~~~~~~ A




% ~~~~~~~~~~~~~~~~~~~~~~~~~~~~~~~~~~~~~~ V frm
\begin{frame}[t,fragile] \frametitle{Palindrome: the idea}
  \begin{columns}[T]
    \begin{column}{0.5\textwidth}
      \small A \textbf{palindrome} reads the same backwards: \texttt{kayak}, \texttt{level},
      \texttt{ey edip adanada pide ye}.

      \vspace{0.4em}
      Compare the two ends, then move inwards:
      \begin{center}
      \begin{tikzpicture}
        \hString[0.55]{k,a,y,a,k}
        \draw[{Stealth}-{Stealth}, thick, orange] (0.27,-0.1) to[bend right=35] (2.47,-0.1);
        \draw[{Stealth}-{Stealth}, thick, blue] (0.82,-0.1) to[bend right=35] (1.92,-0.1);
        \node[font=\scriptsize] at (1.37,-0.95) {\texttt{i} moves right, \texttt{j} moves left};
      \end{tikzpicture}
      \end{center}
    \end{column}
    \begin{column}{0.46\textwidth}
      \SetAlFnt{\small}


% +++++++++++++++++++++++++++++++++++++++ V alg
%: -Alg. 
\begin{algorithm}[H]
  \caption{IsPalindrome($s$)}
  $i \leftarrow 0$, \quad $j \leftarrow \mathit{length}(s) - 1$\;
  \While{$i < j$}{
    \If{$s[i] \neq s[j]$}{
      \Return{false}\;
    }
    $i \leftarrow i + 1$, \quad $j \leftarrow j - 1$\;
  }
  \Return{true}\;
\end{algorithm}
% +++++++++++++++++++++++++++++++++++++++ A alg


    \end{column}
  \end{columns}
\end{frame}
% ~~~~~~~~~~~~~~~~~~~~~~~~~~~~~~~~~~~~~~ A frm




% ~~~~~~~~~~~~~~~~~~~~~~~~~~~~~~~~~~~~~~ V frm
\begin{frame}[t,fragile] \frametitle{Palindrome in Java}
  \lstcoderightstep[0.7][basicstyle=\ttfamily\tiny][basicstyle=\ttfamily\scriptsize]{codeoutput/Palindrome.java}{Java}{codeoutput/Palindrome.txt}

  \vspace{0.3em}
  \small \texttt{toLowerCase()} first, so \texttt{"Level"} counts. The spaces in the last example happen to be symmetric too.
\end{frame}
% ~~~~~~~~~~~~~~~~~~~~~~~~~~~~~~~~~~~~~~ A frm

% ~~~~~~~~~~~~~~~~~~~~~~~~~~~~~~~~~~~~~~ V
\subsection[Count words]{Counting words}
% ~~~~~~~~~~~~~~~~~~~~~~~~~~~~~~~~~~~~~~ A




% ~~~~~~~~~~~~~~~~~~~~~~~~~~~~~~~~~~~~~~ V frm
\begin{frame}[t,fragile] \frametitle{Counting lines, words and characters}
  \lstcoderightstep[0.7][basicstyle=\ttfamily\scriptsize][basicstyle=\ttfamily\tiny]{codeoutput/WordCount.java}{Java}{codeoutput/WordCount.txt}

  \vspace{0.3em}
  \small Input: 3 lines, one empty. \texttt{split("\textbackslash\textbackslash s+")} splits at runs of spaces; \texttt{strip()} avoids an empty first word.
\end{frame}
% ~~~~~~~~~~~~~~~~~~~~~~~~~~~~~~~~~~~~~~ A frm

% ~~~~~~~~~~~~~~~~~~~~~~~~~~~~~~~~~~~~~~ V
\hSection[Summary]{Summary}{}{}
% ~~~~~~~~~~~~~~~~~~~~~~~~~~~~~~~~~~~~~~ A

% ~~~~~~~~~~~~~~~~~~~~~~~~~~~~~~~~~~~~~~ V
\subsection[Rules]{Do and do not}
% ~~~~~~~~~~~~~~~~~~~~~~~~~~~~~~~~~~~~~~ A




% ~~~~~~~~~~~~~~~~~~~~~~~~~~~~~~~~~~~~~~ V frm
\begin{frame}[t,fragile] \frametitle{Do and do not}
  \small
  \begin{columns}[T]
    \begin{column}{0.48\textwidth}
      \begin{block}{Do}
        \begin{itemize}
          \item Count positions from \textbf{0} to \texttt{length()-1}
          \item Loop with \texttt{i < s.length()}
          \item Compare with \texttt{equals}
          \item Store results: \texttt{s = s.trim();}
          \item Put numbers in parentheses: \texttt{"x" + (a + b)}
          \item Call \texttt{nextLine()} after \texttt{nextInt()}
        \end{itemize}
      \end{block}
    \end{column}
    \begin{column}{0.48\textwidth}
      \begin{alertblock}{Do not}
        \begin{itemize}
          \item Compare strings with \texttt{==}
          \item Use \texttt{charAt(s.length())}
          \item Call methods on a \texttt{null} string
          \item Confuse \texttt{'a'} (char) with \texttt{"a"} (String)
          \item Expect \texttt{s.toUpperCase();} alone to change \texttt{s}
        \end{itemize}
      \end{alertblock}
    \end{column}
  \end{columns}
\end{frame}
% ~~~~~~~~~~~~~~~~~~~~~~~~~~~~~~~~~~~~~~ A frm

% ~~~~~~~~~~~~~~~~~~~~~~~~~~~~~~~~~~~~~~ V
\subsection[Exercises]{Exercises}
% ~~~~~~~~~~~~~~~~~~~~~~~~~~~~~~~~~~~~~~ A




% ~~~~~~~~~~~~~~~~~~~~~~~~~~~~~~~~~~~~~~ V frm
\begin{frame}[t,fragile] \frametitle{Exercises}
  \begin{enumerate}\small
    \item Read a full name and print the initials: \texttt{Ayşe Yılmaz} $\to$ \texttt{A.Y.}
    \item Read a word and print it with every letter on its own line, numbered from 1.
    \item Count how many times the letter \texttt{a} (either case) appears in a sentence.
    \item Read a password and say whether it has at least 8 characters, a digit and an upper-case letter
          (use \texttt{Character.isDigit}, \texttt{Character.isUpperCase}).
    \item Read two words and print them in dictionary order (\texttt{compareToIgnoreCase}).
    \item Replace every space in a sentence with \texttt{\_} \textbf{without} \texttt{replace}: build a new string in a loop.
    \item \textbf{Challenge:} check whether a sentence is a palindrome when spaces are ignored
          (\texttt{"Ey Edip, Adana'da pide ye"}).
  \end{enumerate}
\end{frame}
% ~~~~~~~~~~~~~~~~~~~~~~~~~~~~~~~~~~~~~~ A frm




% ~~~~~~~~~~~~~~~~~~~~~~~~~~~~~~~~~~~~~~ V frm
\begin{frame}[t,fragile]
	\begin{center}
		\vspace{4cm}
		\hRemark{\Huge Questions?}
	\end{center}
\end{frame}
% ~~~~~~~~~~~~~~~~~~~~~~~~~~~~~~~~~~~~~~ A frm



\end{document}
